\chapter{Conclusion}
\label{chap:conclusiones}

\lettrine{T}{his} project investigated the use of deep learning techniques for short-term weather forecasting using the ERA5 reanalysis dataset. A forecasting pipeline was implemented to process meteorological data, train neural network models, and evaluate their predictive performance.
The results demonstrate that deep learning models can learn meaningful spatial representations of atmospheric variables and generate realistic temperature forecasts. The U-Net architecture used in this study was able to capture large-scale spatial structures in the temperature field and produce predictions that follow the overall distribution of the observed data. Evaluation metrics such as RMSE and MAE confirmed that the model performs well for short forecast horizons, although prediction errors increase progressively as the forecasting horizon extends further into the future.
Despite these promising results, several limitations remain. The experiments were conducted using only a single meteorological variable (surface temperature), while real atmospheric systems involve complex interactions between many variables such as wind, humidity, and pressure. In addition, the experiments relied on relatively limited computational resources compared to the large-scale infrastructure used in operational forecasting systems.
Future work could therefore focus on integrating additional atmospheric variables, experimenting with more advanced neural network architectures, and conducting larger-scale training experiments using more powerful hardware. Improvements in model design, hyperparameter optimization, and dataset diversity could significantly enhance forecasting accuracy.
In conclusion, this project demonstrates the potential of deep learning approaches for modeling atmospheric dynamics and generating weather forecasts. Although these models are not yet a replacement for traditional numerical weather prediction systems, they represent a promising complementary approach that could contribute to the future development of data-driven weather forecasting methods.

